\documentclass{article}
\usepackage{graphicx} % Required for inserting images

% Added stuff
\usepackage{amsmath}
\usepackage{amssymb}
\newcommand{\N}{\mathbb{N}}
\newcommand{\Z}{\mathbb{Z}}
\newcommand{\Q}{\mathbb{Q}}
\newcommand{\R}{\mathbb{R}}
\newcommand{\C}{\mathbb{C}}
\newcommand{\bigp}[1]{\left(#1\right)}
\newcommand{\bigb}[1]{\left[#1\right]}
\newcommand{\bigabs}[1]{\left|#1\right|}
\newcommand{\summation}[1][0]{\sum_{n=#1}^\infty}
\newcommand{\summationk}[1][0]{\sum_{k=#1}^\infty}
\newcommand{\summationm}[1][0]{\sum_{m=#1}^\infty}

\title{Infinite Series 8}
\author{Siddharth K}
\date{March 2026}

\begin{document}

\maketitle

\section*{Problem 1}
Assume
$$\frac{p}{q} = [a_0; a_1, \ldots].$$
We want to show that it stops after finitely many terms.

Start with:
$$\frac{p}{q} = a_0 + \frac{r_1}{q}.$$
If $r_1 = 0$, the continued fraction has already terminated. If $r_1 > 0$,
$$\frac{p}{q} = a_0 + \frac{1}{q/r_1}.$$
Now divide $q$ by $r_1$:
$$\frac{q}{r_1} = a_1 + \frac{r_2}{r_1},$$
$$\frac{p}{q} = a_0 + \frac{1}{a_1 + \frac{r_2}{r_1}}.$$
If we continue the process,
$$r_{k - 1} = a_kr_k + r_{k + 1}, \quad 0 \le r_{k + 1} < r_k,$$
since each remainder is a nonnegative integer and
$$q > r_1 > r_2 > \cdots \ge 0,$$
the sequence must hit $0$, and the continued fraction stops.

This relates to the Euclidean algorithm because, when computing $\gcd(p, q)$,
\begin{align*}
p &= a_0q + r_1, \\
q &= a_1r_1 + r_2, \\
r_1 &= a_2r_2 + r_3, \\
&\;\;\vdots \\
r_{n - 2} &= a_nr_{n - 1}.
\end{align*}
the quotients $a_0, \ldots, a_n$ are the continued-fraction coefficients.

\section*{Problem 2}
Let $x = [a_0; a_1, \ldots, a_k, \overline{b_1, b_2, \ldots, b_n}]$ and $y = [\overline{b_1: b_2, \ldots, b_n}]$. Let
$$x_i = [a_i; a_{i + 1}, \ldots, a_k, \overline{b_1, b_2, \ldots, b_n}].$$
We claim that each $x_i$ is of the form
$$x_i = \frac{ax + b}{cx + d}$$
for integers $a, b, c, d$.
We can prove this by induction. We have that $x_0 = x$, so the base case holds. Now, assume
$$x_i = \frac{ax + b}{cx + d}.$$
We have that
$$x_{i + 1} = \frac{1}{x_i - a_i} = \frac{1}{\frac{ax + b}{cx + d} - a_i} = \frac{cx + d}{ax + b - a_i(cx + d)} = \frac{cx + d}{(a - a_ic)x + b - a_id}.$$
This means that $x_k$, and therefore $y$, can be written in the form $\frac{ax + b}{cx + d}$. By the same argument as before, $y$ can be written as
$$y = \frac{py + q}{ry + s},$$
for integers $p, q, r, s$. So we have that
$$y(ry + s) = py + q,$$
$$ry^2 + (s - p)y - q = 0.$$
Substituting $y = \frac{ax + b}{cx + d}$,
$$r\bigp{\frac{ax + b}{cx + d}}^2 + (s - p)\bigp{\frac{ax + b}{cx + d}} + q = 0,$$
$$r(ax + b)^2 + (s - p)(ax + b)(cx + d) + q(cx + d)^2 = 0,$$
which is quadratic in $x$ with integer coefficients. So $x$ is a root of a quadratic polynomial with integer coefficients.

\section*{Problem 6}
Let
$$y = \frac{x - 1}{2 + \frac{x - 1}{2 + \cdots}}.$$
Then
$$1 + y = 1 + \frac{x - 1}{2 + y}, \quad y^2 + 2y + 1 - x = 0,$$
$$y = \frac{-2 \pm \sqrt{4 - 4 + 4x}}{2} = -1 \pm \sqrt{x}.$$
Since $y$ is positive,
$$y = -1 + \sqrt{x},$$
$$\sqrt{x} = y + 1,$$
which is the desired continued fraction.

\section*{Problem 7}
Define
$$I_n = \int_0^1 \frac{x^n}{1 + x}dx.$$
We have
$$I_n + I_{n + 1} = \int_0^1 \bigp{\frac{x^n}{1 + x} + \frac{x^{n + 1}}{1 + x}} = \int_0^1 x^n dx = \frac{1}{n + 1},$$
$$\frac{I_{n + 1} + I_n}{I_{n + 2} + I_{n + 1}} = \frac{n + 2}{n + 1}.$$
Define
$$r_n = \frac{I_{n + 1}}{I_n}.$$
Then
$$\frac{1 + 1/r_n}{r_{n + 1} + 1} = \frac{n + 2}{n + 1},$$
$$r_n = \frac{1}{\frac{n + 2}{n + 1}(r_{n + 1} + 1) - 1} = \frac{n + 1}{n + 1 + (n + 2)r_{n + 1}}.$$
So
$$r_0 = \frac{1}{1 + 2r_1} = \frac{1}{1 + \dfrac{2^2}{1 + 3r_2}} = \cdots.$$
We can see that the desired continued fraction is $\frac{1}{1 + r_0}$. Since $r_0 = \frac{I_1}{I_0}$, we have
$$I_0 = \int_0^1 \frac{1}{1 + x}dx = \log(2),$$
$$I_1 = \int_0^1 \frac{x}{1 + x}dx = 1 - \log(2),$$
$$\frac{1}{1 + r_0} = \frac{I_0}{I_0 + I_1} = \log(2).$$

\end{document}